% Part: normal-modal-logic
% Chapter: completeness
% Section: introduction

\documentclass[../../../include/open-logic-section]{subfiles}

\begin{document}

\olfileid{nml}{com}{int}

\olsection{ভূমিকা}

ধরা যাক $\Sigma$ একটি মোডাল পদ্ধতি। বিশুদ্ধতার উপপাদ্য অনুযায়ী, $\Sigma \Proves !A$ হলে $!A$ মডেলের এমন যেকোনো শ্রেণি~$\mClass{C}$-তে বৈধ, যেখানে $\Sigma$-র সব !!{formula}s-এর সব উদাহরণ বৈধ। বিশেষত, $\Log{K} \Proves !A$ হলে $!A$ সব মডেলে সত্য; $\Log{KT} \Proves !A$ হলে $!A$ সব প্রতিবিম্বী মডেলে সত্য; $\Log{KD} \Proves !A$ হলে $!A$ সব সিরিয়াল মডেলে সত্য; ইত্যাদি।

পূর্ণতা বিশুদ্ধতার বিপরীত দিকের দাবি: যেমন, \Log{K} পূর্ণ বলার অর্থ, !!a{formula}~$!A$ বৈধ হলে $\Proves !A$। পূর্ণতা প্রমাণ করা বিশুদ্ধতা প্রমাণের চেয়ে অনেক কঠিন। প্রথমে বিপরীতক্রমটি বিবেচনা করা সুবিধাজনক: \Log{K} তখনই এবং কেবল তখনই পূর্ণ, যখন $\Proves/ !A$ হলেই একটি প্রতিমডেল থাকে, অর্থাৎ এমন মডেল~$\mModel{M}$ থাকে যাতে $\mSat/{M}{!A}$। সমতুল্যভাবে ($!A$-কে নাকচ করে), $\Proves/ \lnot !A$ হলেই $!A$-এর একটি মডেল আছে—এ কথা প্রমাণ করতে পারি। এমন মডেল গড়তে $!A$-তে থাকা তথ্য কাজে লাগাতে পারি। নির্দিষ্ট !!{formula}s-এর মডেল খোঁজার সময়ও আমরা প্রায়ই তা-ই করি: যেমন, $p \lif \Box q$-এর প্রতিমডেল চাইলে জানি, তাতে এমন একটি জগৎ থাকতে হবে যেখানে $p$ সত্য এবং $\Box q$ মিথ্যা। আবার যে জগতে $\Box q$ মিথ্যা, সেখান থেকে অভিগম্য এমন একটি জগৎ থাকতে হবে যেখানে $q$ মিথ্যা। আমাদের জানার জন্য এটুকুই যথেষ্ট: কোন কোন জগতে !!{propositional variable}s সত্য এবং কোন জগৎ থেকে কোন জগৎ অভিগম্য।

কিন্তু পূর্ণতা প্রমাণের ক্ষেত্রে মডেল গড়ার জন্য আগে থেকে কোনো নির্দিষ্ট !!{formula}~$!A$ নেই। আমরা দেখাতে চাই, $\Proves/[\Sigma] \lnot !A$ এমন প্রত্যেকটি $!A$-এর মডেল আছে। এটাই ন্যূনতম দাবি, কারণ \emph{যদি} $\Proves[\Sigma] \lnot !A$ হয়, তবে বিশুদ্ধতা অনুযায়ী $!A$-এর কোনো মডেল নেই (যেখানে $\Sigma$ সত্য)। এবার লক্ষ করি, $\Proves/[\Sigma] \lnot !A$ তখনই এবং কেবল তখনই, যখন $!A$ $\Sigma$-সঙ্গত। (মনে করো, $\Sigma \Proves/[\Sigma] \lnot !A$ এবং $!A \Proves/[\Sigma] \lfalse$ সমতুল্য।) সুতরাং প্রতিটি $\Sigma$-সঙ্গত !!{formula}-এর একটি মডেল গড়াই আমাদের কাজ।

আমাদের কৌশল হলো এমন একটি $\Sigma$-সঙ্গত !!{formula}s-এর সমষ্টি খুঁজে নেওয়া, যাতে~$!A$-এর সঙ্গে আরও সূত্র থাকবে; সেগুলি জানাবে, $!A$-কে সত্য করা জগৎটি কেমন হতে হবে। এগুলি \emph{পূর্ণ} $\Sigma$-সঙ্গত সমষ্টি। $!A$ সত্য করার জন্য একটিমাত্র জগৎ নিয়ে মডেল গড়লেই চলবে না; একাধিক জগৎ এবং তাদের মধ্যে একটি অভিগম্যতা-সম্বন্ধও লাগবে। $!A$-সমন্বিত পূর্ণ $\Sigma$-সঙ্গত সমষ্টিতে $\Box !B$ ও~$\Diamond !C$ আকারের অন্য !!{formula}s-ও থাকবে। প্রতিটি অভিগম্য জগতে $!B$ সত্য হতে হবে, আর অন্তত একটিতে $!C$ সত্য হতে হবে। এজন্য সম্ভাব্য \emph{সব} পূর্ণ $\Sigma$-সঙ্গত সমষ্টিকেই জগৎসমষ্টির ভিত্তি ধরব। মডেলে কোন পূর্ণ $\Sigma$-সঙ্গত সমষ্টি অন্যটি থেকে অভিগম্য হবে, তা নির্ধারণ করাই কঠিন অংশ।

আমরা দেখাব, এভাবে নির্ধারিত মডেলে কোনো জগতে—যেটি নিজেও একটি পূর্ণ $\Sigma$-সঙ্গত সমষ্টি—$!A$ তখনই এবং কেবল তখনই সত্য, যখন $!A$ সেই সমষ্টির !!a{element}। $!A$ $\Sigma$-সঙ্গত হলে সেটি অন্তত একটি পূর্ণ $\Sigma$-সঙ্গত সমষ্টির !!a{element} হবে (এই তথ্যটি আমরা প্রমাণ করব); ফলে এমন একটি জগৎ থাকবে যেখানে $!A$ সত্য। অতএব, আমরা এমন একটিমাত্র মডেল পাব যেখানে প্রত্যেকটি $\Sigma$-সঙ্গত !!{formula}~$!A$ কোনো না কোনো জগতে সত্য। এই মডেলই $\Sigma$-র \emph{প্রামাণিক} মডেল।

\end{document}
